\documentclass[10pt]{article}
% Packages
\usepackage{makecell} % For thicker lines
\usepackage{setspace}  % Controls line spacing
\usepackage{hhline}   % For double horizontal lines
\usepackage{colortbl}  % Add colour to LaTeX tables
\usepackage[T1]{fontenc}  % Choice of font encodings
\usepackage{tgtermes}  % Loads the TeX Gyre Termes font
\usepackage{siunitx}  % A comprehensive (SI) units package
\usepackage{tabularx, booktabs} % For advanced table layout
\usepackage{url}  % Verbatim with URL-sensitive line breaks
\usepackage{authblk}  % For author and affiliation management
\usepackage{natbib}  % A package for bibliographies and citations
\usepackage{graphicx}  % Enhances LaTeX's built-in graphics features
\usepackage{listings}  % Typeset programming code within the document
\usepackage{amssymb}  % Mathematical symbols
\usepackage[nottoc]{tocbibind}  % Adds entries to the table of contents
\usepackage{xcolor}  % Provides easy driver-independent access to colors
\usepackage{microtype}  % Improves the spacing between words and letters
\usepackage{enumitem}  % Control layout of itemize, enumerate, description
\usepackage{tocloft}  % Controls the typographic design of table of contents, etc.
\usepackage[breaklinks,linktocpage]{hyperref}  % Creates hyperlinks in your document
\usepackage[font=small,skip=7pt,labelfont=bf]{caption}  % Customising captions in floating envs

% Adjust the page margins in the bibliography
\let\oldthebibliography=\thebibliography
\let\endoldthebibliography=\endthebibliography
\renewenvironment{thebibliography}[1]{%
  \begin{oldthebibliography}{#1}%
    \raggedright%
    }{%
  \end{oldthebibliography}%
}

\setlength\bibindent{0pt}

% Optional options
% \usepackage{background} % Creates a DRAFT background image on all pages
% \backgroundsetup{contents=Preprint, opacity=0.25, color=gray} % Adds a watermark to the document
% This command changes the line spacing to double.
% ? Needed for reviews/drafts
% \doublespacing

% Custom colours
\definecolor{codegreen}{rgb}{0,0.5,0}
\definecolor{codegray}{rgb}{0.4,0.4,0.4}
\definecolor{codepurple}{rgb}{0.58,0,0.82}
\definecolor{backcolour}{rgb}{0.96,0.96,0.96}
\definecolor{lightgray}{gray}{0.8}

\lstdefinelanguage{JavaScript}{
  keywords={break, case, catch, continue, debugger, default, delete, do, else, finally, for, function, if, in, instanceof, new, return, switch, this, throw, try, typeof, var, void, while, with},
  morecomment=[l]{//},
  morecomment=[s]{/*}{*/},
  morestring=[b]',
  morestring=[b]",
  sensitive=true
}

% Listing styles
\lstdefinestyle{mystyle}{
  frame=tb,
  tabsize=2,
  captionpos=b,
  numbers=left,
  framerule=0pt,
  numbersep=5pt,
  showtabs=false,
  breaklines=true,
  keepspaces=true,
  showspaces=false,
  framextopmargin=6pt,
  framexbottommargin=6pt,
  showstringspaces=false,
  breakatwhitespace=false,
  keywordstyle=\color{purple},
  commentstyle=\color{codegreen},
  stringstyle=\color{codepurple},
  numberstyle=\tiny\color{codegray},
  basicstyle=\ttfamily\footnotesize,
  backgroundcolor=\color{backcolour}}
\lstset{style=mystyle}

% ! Custom template commands
% Add a vertical space after section numbers in ToC
\renewcommand\cftsecafterpnum{\vskip8pt}
% Changes the title of the list of listings
\renewcommand{\lstlistlistingname}{List of \lstlistingname s}
% Changes the title of the bibliography
\renewcommand{\bibname}{References}
% Changes the title of the table of contents
\renewcommand{\contentsname}{Table of Contents}
% Changes the leader between section and page numbers in ToC
\renewcommand{\cftsecleader}{\cftdotfill{\cftdotsep}}
\newcommand{\floatcaption}[2]{\caption[#1.]{#1~#2.}}

% Custom template settings
\captionsetup{justification=centering}  % All captions will be centered
\hypersetup{
  colorlinks = true,
  urlcolor = blue,
  linkcolor = blue,
  citecolor = blue,
  breaklinks = true
}
\PassOptionsToPackage{hyphens}{url}
\urlstyle{same}
\def\Urlmuskip{0mu plus 1mu}
\def\UrlBreaks{\do\/\do-}
\def\UrlBigBreaks{\do\/\do-\do:\do.}
\setlist[itemize]{noitemsep, topsep=0pt, partopsep=0pt}



\begin{document}
\pagenumbering{roman}
\counterwithin{lstlisting}{section}
\counterwithin{figure}{section}
\counterwithin{table}{section}
\setlength{\footskip}{65pt}

% ! ===========================
% ! # MARK: Title, author, etc.
% ! ===========================

\title{\textbf{Death is an \\ Engineering Challenge}}
\author[1]{Daniel Burger}
\affil[1]{\textbf{Eightsix Science Ltd}}
\affil[ ]{\href{mailto:daniel@eightsix.science}{daniel@eightsix.science}}
\author[2]{Masataka Watanabe}
\affil[2]{\textbf{University of Tokyo}}
\affil[ ]{\href{mailto:watanabe@sys.t.u-tokyo.ac.jp}{watanabe@sys.t.u-tokyo.ac.jp}}
\author[3]{Gabriel Cunha}
\affil[3]{\textbf{Tufts University}}
\affil[ ]{\href{mailto:gabriel@neurosyncs.com}{gabriel@neurosyncs.com}}

\author[4]{Izumi Handa}
\affil[4]{\textbf{Panda Lab Inc}}
\affil[ ]{\href{mailto:izumi.handa@pandalab.jp}{izumi.handa@pandalab.jp}}

\date{\textit{\today}}
\maketitle
\thispagestyle{empty}

\begin{sloppypar}

  \begin{figure}[ht]
    \centering
    \includegraphics[width=\textwidth]{figures/cover.png}
    \label{fig:cover}
    % \caption{% TODO: Cover image description.}
  \end{figure}
  \newpage

  % ! ================
  % ! # MARK: Abstract
  % ! ================

  % \begin{abstract}
  %   TODO: Write abstract
  % \end{abstract}

  \pagebreak
  \pagenumbering{Roman}
  \tableofcontents
  \pagebreak
  % \listoffigures
  % \pagebreak
  % \listoftables
  % \pagebreak
  % \addcontentsline{toc}{section}{\lstlistlistingname}
  % \lstlistoflistings
  % \pagebreak
  \pagenumbering{arabic}

  % ! ========================
  % ! # MARK: Document content
  % ! ========================

  \section{Introduction}
  \label{sec:introduction}

  We introduce ‘synconetics’, a new scientific discipline dedicated to solving death through synthetic consciousness mechanics: a set of engineering-focused approaches grounded in solutions already achievable today or in the near future. Synconetics prioritises systems that preserve the physical-dynamical continuity of human consciousness across substrates, sidestepping philosophical speculation in favour of empirical tractability.

  \subsection{First Principles of Death}
  \label{sec:first-principles}

  We frame death not as an inevitability but as a contingent failure of the organised physical system that instantiates consciousness (i.e., the conscious substrate). Physical processes are, in principle, manipulable: no fundamental law prohibits the indefinite persistence of self-maintaining systems under engineered conditions. We, therefore, define dying operationally as the irreversible destruction of the dynamic processes sustaining conscious continuity. By focusing on empirically accessible systems (e.g., the human central nervous system), synconetics circumvents debates about consciousness’s “true nature”—its persistence with what we have at hand becomes an engineering challenge.

  \begin{figure}[ht!]
    \centering
    \includegraphics[width=\textwidth]{figures/4D-process-world-line.png}
    \label{fig:process-world-line}
    \caption{TODO: The 4D process-world-line of a conscious substrate.}
  \end{figure}

  The term conscious substrate denotes any physical medium whose organisation and dynamics have been shown to support conscious experience. Biological substrates, such as the central nervous system (CNS), remain the only empirically confirmed examples. This definition does not restrict consciousness to biology but acknowledges current empirical constraints. Individuals possess direct, if subjective, evidence of their own conscious continuity (e.g., “I think, therefore I am”), providing a provisional anchor for engineering objectives.

  Death marks the irreversible termination of a four-dimensional process world-line, which is the spatiotemporal trajectory of physical states that underpins conscious continuity. Irreversibility reflects thermodynamic reality: entropy increase erodes the recoverability of prior states. Synconetics addresses this through open-system engineering, countering entropy, for example, via energy/matter exchange and error correction. While the precise organisational level critical to consciousness remains unresolved (cellular, molecular, atomic, etc.), engineering pragmatism prevails: we do not need to fully understand a system to work productively with it.

  \subsection{Critiquing Conventional Paradigms}
  \label{sec:new-paradigm}

  Current approaches to overcoming death, whether by delaying or reversing substrate degradation through conventional longevity research (e.g. genetic engineering) or by pursuing substrate independence via computational abstraction (for instance, mind uploading or whole brain emulation), fail to address a fundamental issue: the fragility of the four-dimensional process-world-line that sustains continuous consciousness. Synconetics seeks to safeguard the physical-dynamical continuity of this process while enabling measured transitions to synthetic substrates that promise greater resilience than purely biological ones.

  Mind uploading (MU) and Whole Brain Emulation (WBE) assume that consciousness can survive as a computational abstraction. This is an unverified leap that may pose an existential threat if consciousness proves inseparable from its original physical substrate. Synconetics, therefore, adopts a cautious stance. Synthetic substrates may indeed be more robust, but only with thorough, incremental validation that does not break continuity. By integrating advanced elements such as brain-machine interfaces in ways that preserve the ongoing dynamics of the central nervous system, synconetics aims to maintain an unbroken 4D process-world-line while testing the feasibility of eventually instantiating consciousness in substrates beyond standard biology. Three unresolved challenges exemplify the risks of premature reliance on MU/WBE:

  % TODO Create a figure for each of these
  \begin{enumerate}
    \item \textbf{Philosophical Zombies}: MU/ME’s assumption of substrate independence—that consciousness can “run” on arbitrary hardware—lacks empirical evidence. Computational models simulate neural correlates (e.g., firing rates, connectomes) but cannot verify whether subjective experience persists post-transfer. This leap conflates necessary conditions (physical dynamics like ion gradients and metabolic cycles) with sufficient conditions (unproven for digital substrates). Destructive uploading risks terminating the original process-world-line without verifiable assurance of qualia preservation, effectively gambling one’s existence on unconfirmed substrate equivalency. Even if functional continuity is achieved (à la Chalmers’ ‘gradual replacement’), absent empirical methods to confirm subjective survival, the replica’s consciousness remains an untestable conjecture—a leap no rigorous discipline should endorse.

    \item \textbf{Scale Separability}: Emulating brains at biologically relevant scales demands energy exceeding Earth’s projected budget for centuries. Modern supercomputers simulate <1\% of a human brain’s synaptic activity at 1,000\(\times\) slowed time—a trivial fraction of real-time biological efficiency (20W vs. megawatts). Crucially, the brain’s dynamics blend deterministic chaos and stochastic noise as quantum-scale fluctuations (e.g., ion channel gating) propagate into macroscopic neural activity. This “web of causality” \citep{watanabe_biological_2022}—though shorter-lived than classical models suggest—defies reduction to deterministic simulations. Worse, fidelity requirements escalate exponentially if consciousness depends on quantum effects (e.g., microtubule coherence as proposed by Orch-OR). To simulate reality, we asymptotically approach infinity by simulating it within itself (i.e. infinite recursion).

    \item \textbf{Teleportation Paradox}: Most MU/WBE approaches rely on destructive scanning or non-destructive copying to transfer consciousness onto a computational substrate. Both methods sever the original’s causal continuity: destructive scans terminate the 4D process-world-line outright, while non-destructive copies spawn a parallel entity divorced from your subjective stream. Even “perfect” replication creates a new 4D trajectory; you do not experience the replica’s existence. Critics dismiss this as philosophical nitpicking—“I don’t care if it’s a copy, as long as it thinks it’s me.” But survival hinges on \emph{your} conscious continuity, not a replica’s beliefs. Even if half of humanity accepted copy-based “survival,” the other half would reject existential roulette. Synconetics prioritises solutions that preserve the 4D process-world-line outright—ensuring no one is forced to gamble their existence on untested metaphysics.
  \end{enumerate}

  Critics accusing biological bias misunderstand the burden of proof: synthetic substrates must demonstrate conscious continuity before replacing biology. Until then, privileging verified systems is engineering prudence, not chauvinism. The imperative is clear: abandon replication metaphysics and address substrate fragility of the 4D process-world-line—the root cause of death.

  \subsection{Synconetics: Establishing a New Discipline}
  \label{sec:new-discipline}

  The inadequacy of conventional paradigms demands more than incremental refinement—it necessitates a foundational realignment of objectives and methods. We formally establish synconetics as a distinct discipline focused on achieving synthetic consciousness mechanics (SCM): the applied science of sustaining conscious continuity through substrate stabilisation, repair, and non-destructive transition. This demarcation is structurally essential, not terminological, arising from three irreconcilable divergences from prior approaches:

  \begin{enumerate}
    \item \textbf{Process-World-Line Fidelity}: Exclusive focus on preserving the unbroken 4D causal chain of the conscious substrate’s physical dynamics.
    \item \textbf{Non-Negotiable Continuity}: Rejection of destructive replication (copy/transfer, scanning/emulation) as existential threats, not technical shortcuts.
    \item \textbf{Phenomenological Grounding}: Subjective reportability during interventions as the gold standard for success, overriding third-person behavioural metrics.
  \end{enumerate}

  Synconetics is defined by its engineering pragmatism: it employs only technologies feasible today—closed-loop BCIs, in vivo neuroprosthetics, and metabolic stabilisation systems—to extend substrate viability while developing transitional synthetic components. Crucially, it rejects speculative futurism (e.g., “post-biological consciousness”), lacking mechanistic pathways from current capabilities.

  Ethically, Synconetics introduces a Non-Destruction Principle (NDP): any intervention must preserve the original substrate’s causal continuity unless synthetic components achieve parity under Husserlian bracketing (i.e., indistinguishable subjective experience during phased integration). This principle reframes existential risk mitigation as an engineering constraint—a system that cannot prove continuity preservation during testing is axiomatically unfit for deployment. Contrast this with MU/WBE’s ethical vacuum, where the destruction of the biological original is often a prerequisite.

  Critics may challenge Synconetics’ conservative materiality—its insistence on privileging biological substrates until synthetic alternatives meet NDP standards. We counter that this stance is empirically enforced:


  \begin{itemize}
    \item P1: All empirically verified consciousness resides in biological substrates.
    \item P2: Substrate transitions risk discontinuity until synthetic alternatives are validated.
    \item C: Therefore, biological conservation is the null hypothesis; synthetic integration requires affirmative proof.
  \end{itemize}

  Verification protocols operationalise this:

  \begin{enumerate}
    \item Real-time neural-phenomenological correlation during interventions (e.g., EEG signatures of self-awareness concurrent with first-person reporting).
    \item Thermodynamic parity checks ensuring synthetic components operate within biological free energy ranges.
    \item Plasticity synchronisation metrics confirming synthetic-biological interplay respect natural neural timescales (e.g., dopamine-based reinforcement cycles).
  \end{enumerate}

  Synconetics’ immediate priorities reflect its engineering focus:

  \begin{itemize}
    \item Substrate Stabilisation: Deploying existing technologies—nanoparticle oxygen carriers, targeted hypothermia—to extend the “recoverability window” post-clinical death from minutes to hours.
    \item Continuity-Assured Augmentation: Gradual CNS hybridisation using neural prosthetics validated for dynamical equivalence (e.g., replacing cerebellar circuits with neuromorphic chips only after in situ testing confirms motor learning rates remain physiologically consistent).
    \item Failure Mode Cataloguing: Systematically mapping critical collapse pathways (e.g., glutamate storms, microtubule decoherence) to prioritise intervention targets.
  \end{itemize}

  Synconetics’ founding axiom—that death is an engineering challenge, not a metaphysical inevitability—demands this disciplinary independence. Only through dedicated convergence of its constituent fields can we systematically dismantle the technical failures underlying conscious collapse. To treat it as a subfield of bioengineering or computer science would be to repeat the failures of “AI safety” as an afterthought to capability development. Here, continuity is the capability—and its assurance requires a discipline born of that singular purpose.

  \subsection{Foundational Principles of Synconetics}
  \label{sec:foundational-principles}

  Synconetics’ methodology is governed by five non-negotiable axioms that constrain both its engineering objectives and ethical boundaries. These principles operationalise the discipline’s core innovations, transcending mere philosophical stances to define actionable design constraints.

  % TODO Have a table where we differentiate between the principles and the conventional paradigms.

  Principle 1: Contingency of Death as Substrate Failure Mode
  Death is reclassified as a systems engineering problem—specifically, the failure to maintain critical dynamics in a thermodynamically open conscious substrate. This reframing rejects metaphysical fatalism, focusing instead on three tractable engineering targets:
  - Failure Mode Analysis: Cataloguing collapse pathways (e.g., ischemic cascades, protein misfolding avalanches) as prioritised intervention points.
  - Recoverability Window Extension: Using existing clinical technologies (e.g., cryoprotective perfusion, nanoparticle O2 carriers) to push the post-clinical-death reversibility horizon from minutes to hours.
  - Resilience by Design: Architecting synthetic-biological hybrids to tolerate subsystem failures without process-world-line termination (e.g., distributed neuromorphic backups for hippocampal CA3 circuits).

  Principle 2: Engineering First, Speculation Last
  Methodological primacy is given to physically instantiable solutions using 2020s-era technologies, with theoretical commitments minimised to:
  - Weak Physicalism: If consciousness depends on substrate dynamics, sustaining those dynamics sustains consciousness (falsifiable via first-person discontinuities during intervention).
  - Non-Magic Axiom: No unproven physics (e.g., quantum consciousness) are invoked unless they provide testable engineering advantages.

  This principle bans common distractions:
  - No resources allocated to “consciousness detectors” for synthetic systems—biological continuity provides the only verification needed.
  - No engagement with theory wars (IIT vs. Global Workspace)—instead, targeting conserved neural correlates like thalamocortical resonance.

  Principle 3: Continuity as Causal Topology, Not Mere Persistence
  Continuity is defined not as indefinite substrate survival but as preserving the autopoietic causal graph—the self-sustaining web of microphysical interactions (synaptic, metabolic, electrochemical) that constitute the process-world-line. Key engineering implications:
  - Intervention Threshold: Any modification must retain >99.9\% of causal interactions per millisecond (derived from neural network error-correction limits).
  - Non-Locality Constraints: Components replaced must not participate in cross-scale dynamics (e.g., microtubule quantum coherence) until synthetic analogues achieve femtosecond-scale bioalignment.

  Principle 4: Substrate Plasticity as the Transition Engine
  Physical realisability is achieved not through brute-force emulation but by leveraging the CNS’s innate adaptive capacity to integrate synthetic components. This mirrors cochlear implant neurointegration but with stricter continuity safeguards:
  - Plasticity Alignment Protocol: Synthetic interfaces must operate within biological learning rate bounds (e.g., dopamine-driven reinforcement windows of 50-500ms).
  - Phase-Locked Replacement: Gradual substitution of neural circuits only during their native quiescent periods (e.g., replacing hippocampal place cells during slow-wave sleep cycles).

  Principle 5: Survival Trumps All Optimisation Pressures
  Enhancements (e.g., expanded working memory via neuromodulation) are permissible only if:
  - Zero Discontinuity Risk: Enhancement provides survival redundancy (e.g., memory augmentation doubles as stroke recovery failsafe).
  - No Emergent Goal Conflicts: Augmentation cannot create incentive structures that undervalue continuity (e.g., cognition-speed boosts that destabilise emotional salience networks).

  Ethical enforcement mechanisms:
  - Enhancement Moratorium Clause: Any intervention with >0.1\% discontinuity risk in preclinical models is banned, regardless of ancillary benefits.
  - Survival-Benefit Ratio (SBR): Enhancements must demonstrate a minimum 10:1 survival-to-enhancement resource allocation.

  These principles collectively prevent mission drift into transhumanist abstraction. By binding Synconetics to the causal physics of existing conscious systems, they ensure the discipline remains an engineering—not philosophical—endeavour.

  \section{Nomenclature and Definitions}
  \label{sec:nomenclature}

  Establishing Synconetics as a rigorous discipline demands terminological precision to avoid conceptual drift endemic to consciousness studies. Below, we define core concepts, contrasting them with conventional paradigms while maintaining strict alignment with the foundational principles outlined in Section~\ref{sec:introduction}.

  \subsection{Core Definitions}
  \textbf{Synconetics}: The engineering discipline dedicated to sustaining conscious continuity by treating biological death as a substrate failure mode. It combines neuroengineering, systems thermodynamics, and phenomenological validation to extend the 4D process-world-line of consciousness-supporting substrates. Unlike speculative mind-uploading paradigms, Synconetics prioritises \textit{physical persistence} over abstract replication.

  \textbf{Synthetic Consciousness Mechanics (SCM)}: Applied methodologies for designing, interfacing, and transitioning consciousness-supporting substrates. SCM’s scope includes:
  \begin{itemize}
    \item Substrate stabilisation (e.g., cryoprotective perfusion protocols)
    \item Non-destructive integration of synthetic components (e.g., plasticity-aligned neuroprosthetics)
    \item Real-time causal continuity monitoring (see Principle 3, Section~\ref{sec:foundational-principles})
  \end{itemize}

  \textbf{Synthetic Consciousness Substrate (SCS)}: Any physical system—biological, synthetic, or hybrid—capable of sustaining the dynamic processes underlying conscious experience. An SCS must:
  \begin{itemize}
    \item Operate as a thermodynamically open system, exchanging energy/matter to maintain non-equilibrium dynamics
    \item Exhibit neural-phenomenological isomorphism with biological benchmarks (e.g., thalamocortical resonance at 40 Hz gamma synchrony)
    \item Support millisecond-scale feedback loops essential for metacognition
  \end{itemize}
  \textit{Note}: Biological brains remain the sole empirically validated SCS; synthetic variants require validation via first-person reportability during phased integration.

  \textbf{Synthetic Consciousness Transfer (SCT)}: Protocols enabling migration of conscious processes between substrates \textit{without} interrupting the causal topology of the 4D process-world-line. SCT expressly prohibits destructive scanning/emulation, instead employing:
  \begin{itemize}
    \item Gradual neuroprosthetic replacement validated by real-time neural-phenomenological correlation
    \item Closed-loop BCI mediation ensuring <0.1\% deviation in critical dynamics (e.g., spike-phase alignment in hippocampal theta cycles)
  \end{itemize}
  \textbf{Continuity} is mathematically defined as:
  \begin{equation}
    C(t) = \int_{t_0}^{t_1} \frac{\partial S}{\partial t} \, dt \quad \text{where } \Delta S_{\text{critical}} < \epsilon_{\text{plasticity}}
  \end{equation}
  Here, \( S \) represents substrate state vectors, and \( \epsilon_{\text{plasticity}} \) denotes the maximum allowable deviation before neural adaptive mechanisms fail (derived from spike-timing-dependent plasticity thresholds).

  \textbf{Synthetic Consciousness Interfacing (SCI)}: Bidirectional systems coupling an SCS to external environments while preserving agency and phenomenological coherence. SCI must achieve:
  \begin{itemize}
    \item Latency <20 ms for sensorimotor loops (matching biological visuomotor integration)
    \item Entropic parity between synthetic and biological perceptual inputs.
  \end{itemize}

  \textbf{Process-World-Line}: The unique spatiotemporal trajectory of a consciousness-supporting substrate’s dynamics. Collapse of this trajectory—through irrecoverable state loss or causal discontinuity—constitutes death under Synconetics’ framework.

  \subsection{Terminological Distinctions}
  \begin{table}[ht]
    \centering
    \caption{Contrasting Synconetics Terminology with Conventional Paradigms}
    \begin{tabular}{p{0.3\textwidth}p{0.3\textwidth}p{0.3\textwidth}}
      \toprule
      \textbf{Synconetics Term}               & \textbf{Conventional Counterpart} & \textbf{Key Distinction}                                        \\
      \midrule
      Synthetic Consciousness Transfer (SCT)  & Mind Uploading                    & Maintains causal continuity; rejects destructive scanning       \\
      Synthetic Consciousness Substrate (SCS) & Computational Emulation           & Requires thermodynamic openness; not limited to digital systems \\
      Continuity (Processual)                 & Pattern Identity                  & Defined by physical dynamics, not informational isomorphism     \\
      Substrate Stabilisation                 & Lifespan Extension                & Targets process-world-line resilience, not biological longevity \\
      \bottomrule
    \end{tabular}
    \label{tab:terminology}
  \end{table}

  \textbf{Mind Uploading (MU)}: Often implies destructive scanning followed by computational replication. Synconetics classifies MU as \textit{replica creation}, not survival, due to causal discontinuity.

  \textbf{Artificial Consciousness}: Typically denotes de novo creation of conscious systems. Synconetics focuses exclusively on \textit{continuity of existing individuals}, making this term inapplicable.

  \textbf{Machine Consciousness}: Assumes implementation on classical computing architectures. Synconetics rejects this substrate chauvinism, as SCSs may exploit quantum biological or neuromorphic properties.

  \subsection{Resolving Ambiguities}
  To prevent inconsistent usage:
  \begin{itemize}
    \item \textbf{Substrate}: Always denotes \textit{consciousness-supporting} systems. For non-conscious components (e.g., robotic limbs), use \textit{peripheral module}.
    \item \textbf{Transfer (SCT)}: Never used without “Synthetic Consciousness” prefix to avoid conflating with data transmission.
    \item \textbf{Process-World-Line}: Hyphenated consistently, reflecting its 4D spacetime ontology (cf. Section~\ref{sec:first-principles}).
  \end{itemize}

  This lexicon provides the semantic rigour necessary to engineer death’s defeat without metaphysical detours. By binding terminology to physically measurable quantities and testable interventions, Synconetics transcends the equivocation that has long plagued consciousness research.

  % TODO Non-fixed yet

  \section{Feasibility and Opportunities}
  \label{sec:feasibility}

  These examples expose the substrate’s elusive boundaries: consciousness likely depends on dynamics within the CNS, not isolated structures. This empirical approach drives two interdependent strategies:

  \begin{itemize}
    \item \textbf{Decoupling agency from embodiment}: Biological embodiment conflates the substrate (brain) with its environmental interface (body), creating existential fragility. As an example, Modern computing systems avoid single points of failure by decoupling data from hardware: critical servers reside in shielded facilities with redundant power and backups, surviving even if individual nodes are destroyed. Synconetics applies this principle to consciousness. A brain interfacing with a virtual avatar or robotic form—untethered from its biological casing—shields the substrate in controlled environments, while agency persists via peripheral embodiments. Death becomes contingent not on localised physical destruction but on breaching the substrate’s engineered fault tolerance.
    \item \textbf{Redundancy through resolution}: Humans repair computers with near-atomic precision yet lack tools to correct even minor neural degradation—an indefensible asymmetry. High-bandwidth neural interfaces resolve this absurdity, enabling synthetic augmentation while preserving the 4D process-world-line. Higher resolution permits repairs once deemed impossible, transforming the brain from a “black box” to an engineerable substrate. <also say that this allows augmentation as we have the full control of consciousness>
  \end{itemize}

  The trajectory is unambiguous. First, preserve the brain’s native dynamics. Next, harden it through synthetic augmentation—each intervention reducing the system’s thermodynamic “surface area” for collapse. Finally, achieve full substrate independence, rendering death not just solvable but physically implausible, whether the substrate persists as biological tissue, biohybrid systems, or synthetic materials.

  A core assertion of this essay is that addressing the cessation of consciousness—biological death—is not merely a future aspiration contingent upon resolving the deepest mysteries of mind, but a challenge amenable to engineering methodologies today. The establishment of Synconetics rests upon the conviction that a practical, near-term research and development programme is feasible, grounded in established scientific principles and leveraging current technological trajectories. This section outlines the principles underpinning this feasibility.

  History repeatedly demonstrates that transformative engineering often precedes complete scientific elucidation; thermodynamics was harnessed before statistical mechanics provided a full explanation, and controlled flight was achieved before fluid dynamics was comprehensively understood. Synconetics adopts a similar pragmatic posture. While acknowledging the profound complexity of consciousness, potentially exceeding historical analogies, our approach focuses on manipulating and interfacing with the known physical substrate using established physical and biological principles. It does not predicate itself on first achieving a final, universally accepted theory of consciousness—a pursuit that, while valuable, may be indefinitely protracted. Instead, Synconetics targets the engineering problem of preventing the failure of the system currently known to support consciousness.

  Crucially, Synconetics does not require the invention of entirely new scientific fields ex nihilo. Its feasibility is substantially bolstered by building directly upon the rapid, ongoing advancements occurring across a range of synergistic domains. Progress in neuroscience yields increasingly detailed maps of neural circuits, deeper understanding of plasticity, and refined identification of the neural correlates of consciousness. Neuroengineering provides increasingly sophisticated tools for brain-computer interfacing (BCIs), neurostimulation, and high-resolution neural recording. Materials science offers novel biocompatible and ‘smart’ materials essential for interfacing and substrate construction. Bioengineering contributes techniques in tissue engineering, organoid development, and stem cell technology. Robotics and AI provide sophisticated control systems and simulation environments vital for interfacing and testing. Synconetics’ viability arises from the strategic integration and focused application of these converging capabilities towards the specific goal of ensuring conscious continuity.

  Methodologically, Synconetics derives strength from its focus on physical processes rather than abstract definitions. By concentrating on the tangible substrate and its continuous dynamic activities, the problem is framed in terms of measurable, manipulable physical variables. This approach circumvents the immediate need to resolve intractable philosophical debates about the essential nature of mind, information versus matter, or the precise definition of qualia—debates that often stall progress in paradigms reliant on abstract computationalism or functional equivalence alone. The engineering target becomes the verifiable preservation of the physical process known to underpin consciousness in the individual. This focus on the physical is, admittedly, a working hypothesis—it assumes a sufficiently tight coupling between the targeted physical dynamics and the subjective experience they support. However, from an engineering perspective aimed at preserving an existing conscious system, maintaining the integrity of its known physical basis represents the most rational, conservative, and empirically grounded strategy currently available. It is an approach rooted in physicalism and prioritises non-destruction of the only confirmed instance of the phenomenon we seek to preserve.

  The brain’s inherent plasticity—its remarkable capacity for functional and structural reorganisation in response to learning, injury, or environmental changes—provides another key enabling factor. Synconetics methodologies, particularly those involving gradual intervention such as progressive repair, augmentation, or substrate replacement (as exemplified in Section \ref{sec:daniel-approach}), are designed to leverage this natural adaptability. The hypothesis, supported by neurological precedent, is that gradual, carefully managed changes can allow the neural system to adapt and maintain functional and informational continuity throughout the transition, mitigating the profound risks associated with abrupt, large-scale alterations. Determining the precise limits of plasticity, particularly concerning the faithful preservation of identity-critical information like specific memories or personality traits during substantial structural change, remains a crucial area of research within Synconetics.

  Furthermore, achieving meaningful interaction with an environment—a critical aspect of Synthetic Consciousness Interfacing (SCI)—may not necessitate the perfect replication of biological sensory fidelity. Existing sensory prosthetics, such as cochlear and retinal implants, alongside increasingly sophisticated BCIs, demonstrate that functionally sufficient interaction with external environments (whether physical or virtual) can be achieved, even if the subjective quality of the experience is altered or simplified. This suggests the interfacing challenge, while significant, might be more tractable than achieving full biological equivalence, focusing instead on providing the necessary bandwidth and control for agency and coherent experience. Defining ‘sufficiency’ in this context, particularly regarding long-term psychological well-being and the preservation of a rich sense of self, remains an important ongoing challenge requiring input from phenomenology and psychology.

  The overall engineering challenge, though immense, also appears decomposable into distinct, albeit interconnected, sub-problems. These include the development and validation of alternative consciousness-supporting substrates (SCS), the devising of reliable, continuity-preserving transfer or transition methodologies (SCT), and the creation of effective, high-bandwidth interfaces (SCI). This potential modularity allows for focused research and development efforts within specialised teams, mirroring standard practice in complex engineering projects and making the overarching goal seem less monolithic and more approachable through parallel advancements.

  Synconetics also gains feasibility by deliberately seeking to avoid reliance on distant or highly speculative scientific breakthroughs. Unlike paradigms potentially dependent on future revolutions in fundamental physics, practical quantum computing, the emergence of artificial general intelligence for validation, or atomically precise nanotechnology for scanning and construction, Synconetics aims to progress primarily by pushing the boundaries of technologies and scientific principles that are either available now or represent foreseeable extensions of current capabilities. While significant advancements are undoubtedly required—for instance, in long-term stable bio-hybrid materials, minimally invasive large-scale neurosurgery, or reliable ex vivo organ support—these largely fall within the projected trajectory of contemporary bioengineering, materials science, and neurotechnology. The definition of ‘foreseeable’ remains inherently subjective, yet the principle guides Synconetics towards solutions grounded in known physics and biology.

  Finally, the pursuit of Synconetics’ long-term objectives is expected to yield valuable scientific knowledge and intermediate technological applications. Advancements in neural modelling, regenerative therapies for brain injury and neurodegenerative diseases, radically improved BCIs for communication and control, and novel biocompatible materials are all likely spin-offs. This potential for generating near-term scientific, therapeutic, and potentially commercial value provides a pragmatic justification for investment and effort, offering tangible returns even before the ultimate goal of indefinite conscious continuity is achieved. This aligns Synconetics with a model of progressive innovation, where intermediate milestones contribute significantly in their own right.

  % ! MARK: Methodologies and Approaches
  \section{Methods and Approaches}
  \label{sec:methods}

  Synconetics, having been delineated through its foundational principles and core engineering focus, demands translation into practical methodologies capable of achieving its stated objectives. A essay without demonstrable avenues for realisation remains purely theoretical. To affirm that this framework guides tangible research and development efforts commencing today, this section presents two distinct approaches currently under investigation by the authors. These serve as initial exemplars of Synthetic Consciousness Mechanics in practice, illustrating how the formidable challenge of preserving conscious continuity can be addressed through concrete, verifiable engineering strategies. Each methodology aligns with the core Synconetics tenets—particularly the prioritisation of physical process continuity—and leverages contemporary neuroscience and engineering capabilities. While representing only the first steps within this nascent discipline, they stand distinct from paradigms reliant on destructive replication or purely abstract computation, showcasing the pragmatic potential inherent in the Synconetics approach.

  \subsection{Ectopic Cognitive Preservation}
  \label{sec:daniel-approach}

  The strategy termed ‘Ectopic Cognitive Preservation’ (ECP), under development by Eightsix Science\footnote{Disclosure: Daniel Burger is a co-founder of Eightsix Science.}, exemplifies a Synconetics methodology squarely focused on ensuring the physical continuity of the biological substrate through gradual, technologically mediated replacement. Its core technical proposal involves the progressive, piecemeal substitution of existing biological brain tissue with bio-hybrid neural grafts. These grafts are envisaged as constructs of living neural tissue, derived from the patient’s own induced pluripotent stem cells (autologous iPSCs) differentiated into appropriate neural lineages to circumvent immune rejection, and integrated during advanced bioprinting (potentially 4D bio-hybrid printing) with micro- or nano-scale electronic components. These integrated elements could serve various functions, such as sensing local activity, providing targeted stimulation, offering structural support, or facilitating metabolic exchange. Achieving and rigorously verifying true functional equivalence between the original tissue and the graft—encompassing not merely basic neuronal firing but complex network dynamics, synaptic plasticity profiles, and the preservation of identity-critical information patterns—represents a monumental, yet central, challenge for this approach.

  ECP’s commitment to continuity hinges critically on the principle of gradualism, designed to leverage the brain’s inherent plasticity and capacity for functional reorganisation, analogous to adaptations observed in response to slow-growing lesions like benign gliomas. The core hypothesis is that by carefully managing the rate of replacement—gradually silencing small portions of original tissue while simultaneously integrating new, active bio-hybrid grafts—neural information processing and functional roles can migrate or be re-encoded within the new substrate without disrupting the overall continuity of cognitive processes and, crucially, conscious experience. This reliance on plasticity, while biologically plausible, carries inherent risks regarding the fidelity of information preservation; ensuring that specific memories, learned skills, and personality nuances are faithfully maintained during such transitions, rather than merely enabling functional adaptation, remains a key area requiring deep theoretical understanding and empirical validation. Methodologies for precisely controlling gradual silencing and for real-time monitoring of graft integration and functional takeover are therefore critical research components.

  The initial outcome targeted by ECP is a rejuvenated, potentially enhanced biological or bio-hybrid brain residing within the original cranium. Composed progressively of the new graft material integrated with embedded electronics, this enhanced substrate aims primarily to halt or reverse age-related degradation within the brain itself, thereby addressing a primary failure mode of the current biological system. The integrated electronics could also offer inherent capabilities for advanced Brain-Computer Interfacing (BCI), potentially enabling seamless integration with virtual or augmented reality environments without requiring separate invasive procedures later.

  The ultimate, more radical goal of ECP involves the surgical explantation of this fully replaced bio-hybrid brain. Sustained long-term via an advanced, closed-loop whole-brain perfusion system providing a meticulously controlled physiological environment ex vivo, its function would be embedded within sophisticated virtual environments through high-bandwidth communication channels derived from the integrated electronics. This step aims to achieve complete decoupling from the vulnerabilities of the original biological body. Realising stable, long-term ex vivo maintenance presents immense technical hurdles, demanding perfect replication of complex physiological conditions. Furthermore, the profound ethical and psychological implications of explantation and existence within a potentially constrained virtual reality necessitate careful consideration far beyond mere technical feasibility, touching upon questions of identity, well-being, and the nature of experience itself, reminiscent of speculative explorations like the “San Junipero” thought experiment but demanding rigorous, real-world analysis.

  Methodologically, ECP aligns directly with the foundational principles of Synconetics. It treats death as a substrate failure problem (Principle 1), employs a tangible engineering methodology involving bioprinting, grafting, BCI, and perfusion systems (Principle 2), and explicitly prioritises physical process continuity through gradual replacement, rejecting destructive copying (Principle 3). Its feasibility is grounded in known physical and biological processes like cell differentiation, neural plasticity, and electronics integration (Principle 4), and its initial focus is squarely on survival and resilience by halting ageing and enabling repair, subordinating potential enhancements (Principle 5). While the projected timelines and the ultimate certainty of guaranteeing continuity face valid scrutiny and require significant empirical validation, ECP serves as a concrete example of the Synconetics paradigm: pursuing ambitious engineering goals within a framework of physical continuity and ethical caution. Furthermore, its roadmap inherently generates intermediate technologies—advanced neural simulation, high-fidelity graft production, progressive replacement techniques—with significant near-term therapeutic and research value (e.g., for neurodegenerative disease, personalised medicine), offering a pragmatic pathway for development and funding.

  \subsection{“Masataka’s Approach”}
  \label{sec:masataka-approach}

  % TODO Say somewhere that this is not classical mind uploading but something new!
  \begin{itemize}
    \item To avoid death, one must upload their consciousness while still alive. Studies on split-brain patients show that separating the brain’s hemispheres results in two consciousnesses, while reconnecting them fuses them back into one. This insight leads to the first step: separating the hemispheres and linking each to an mechanical one using a brain-machine interface (BMI).
    \item Once connected, consciousness and memory begin integrating across biological and mechanical hemispheres. Just like patients who lose one hemisphere, consciousness can persist entirely in the remaining half. Likewise, when the biological hemisphere dies, consciousness continues in the mechanical one.
    \item To complete the process, the two mechanical hemispheres—each paired with a biological one—are reconnected, restoring unity. Thus, the transition to an mechanical brain is seamless, without death.
    \item The process begins with a craniotomy to insert a segmented BMI between hemispheres. This splits your consciousness into two—each aware only of one side of your body and world. Though disorienting, it’s temporary and safe.
    \item The mechanical hemisphere initially holds no personality. Through neural routing, connections form between the hemispheres. Gradually, lost sensations and movements return. With full routing, the two halves function as one.
    \item Memory transfer follows. In the first session, recalling memories stores them in the mechanical hemisphere. In the second, electrical stimulation triggers forgotten memories, transferred similarly unless blocked by the user. Painful memories may be skipped, though they often shape identity.
    \item Finally, as the biological hemisphere fails, you lose right-side perception but retain full awareness via the mechanical side. Then, both mechanical hemispheres are merged, completing the upload. Consciousness lives on digitally, potentially returning to the physical world as an avatar.
    \item Actual Procedures
    \item Welcome to the World of Split Brains: Insertion of the Brain-Machine Interface
    \item To begin the consciousness upload, we implant a segmented brain-machine interface (BMI) via a craniotomy, placing it between the major fiber bundles connecting the brain’s hemispheres. This is similar to a corpus callosotomy—a routine epilepsy treatment—and is considered safe.
    \item However, the BMI will split your consciousness into two entities, like in split-brain patients. For now, imagine you’re in the left hemisphere—you’ll perceive, feel, and control only the right side. This disconnection is temporary and necessary on the path to digital immortality.
    \item Connection with the Mechanical Hemisphere. The BMI is coated with proteins that help neurons form new synapses. After a few days, it’s ready to connect with an mechanical hemisphere that starts off with no memories or personality.
    \item Initially, you may notice faint sensations on the left side of your body or visual field. As “neural routing” begins, neurons from your biological hemisphere connect to the mechanical one. Vision, movement, and bodily awareness slowly return, starting off distorted and becoming more clear and coherent. Eventually, both hemispheres function as one, fully integrating consciousness.
    \item Memory transfer begins with two sessions. In the first, you actively recall memories, which are mirrored in the mechanical hemisphere and stored long-term—without conscious notice.
    \item In the second, a flashback mechanism triggers long-lost memories through electrical stimulation, based on Penfield’s experiments. Vivid past scenes resurface and are recorded unless you press a button to stop painful memories from transferring. These can shape identity, but suppressing all may dilute character. Even after uploading, forgotten memories might return unexpectedly. This can be controlled by weakening related synaptic links in the mechanical brain.
    \item The Final Phase: Termination of the Biological Hemisphere Once the process is complete, your consciousness spans your biological left and mechanical right hemispheres. When your biological brain shuts down, sensations from the right side vanish—but your consciousness persists in the mechanical hemisphere.
    \item When the biological hemisphere dies, consciousness doesn’t. Like surviving a stroke that wipes out one side of your brain, your mind shifts seamlessly into the mechanical hemisphere. Death is bypassed; your self continues uninterrupted.
    \item The final step merges both mechanical hemispheres, previously linked to each side of your brain. Now unified, your consciousness lives fully in a digital substrate.
    \item Real-time operation in the physical world is costly, so your mind transitions to a virtual world with others who’ve uploaded. Someday, you might return as an avatar—if you can pay the fee.
  \end{itemize}


  % ! MARK: A Pragmatic Roadmap for Synconetics
  \section{Roadmap}
  \label{sec:roadmap}

  The establishment of Synconetics as a viable discipline demands more than foundational principles and precise nomenclature; it requires a pragmatic research and development roadmap. Such a roadmap must candidly acknowledge the profound technical and conceptual challenges inherent in engineering conscious continuity, whilst simultaneously identifying tractable starting points and strategic pathways forward. A central tenet of this pragmatism is the parallel pursuit of complementary methodologies. These distinct approaches address different facets of the core problem, mitigate different categories of risk, and strategically leverage both current and foreseeable technological capabilities. This section outlines such a strategic approach, illustrating the engineering feasibility central to Synconetics by focusing on the interplay between biologically grounded interventions and the exploration of synthetic substrates.

  Given the persistent uncertainties surrounding the precise physical prerequisites for consciousness and the optimal characteristics of a long-term, resilient substrate, a prudent strategy necessitates pursuing distinct, yet potentially synergistic, research programmes concurrently. We advocate for a dual-pronged approach. The first prong encompasses methodologies focused on preserving continuity by directly augmenting, repairing, or gradually replacing the existing biological substrate. This biologically grounded path, exemplified by the Ectopic Cognitive Preservation (ECP) strategy detailed in Section \ref{sec:daniel-approach}, prioritises working with the known substrate, leveraging established biological mechanisms such as neural plasticity alongside advancements in tissue engineering, bio-hybrid integration, and neurosurgery. Its initial focus lies squarely on mitigating intrinsic biological failure modes, primarily age-related degeneration, and enhancing the resilience of the existing system. While potentially offering a nearer-term route by sidestepping the challenge of creating consciousness ex nihilo in an entirely novel medium, this approach ultimately retains a substrate with inherent biological vulnerabilities. Even an advanced bio-hybrid brain, particularly if maintained ex cranio, remains susceptible to physical destruction, lacks intrinsic fault tolerance compared to potentially achievable engineered systems, and may face fundamental biological limitations that constrain indefinite persistence.

  The second prong of our strategy directly confronts the challenge of engineering non-biological or radically different physical systems capable of supporting conscious processes, coupled with developing rigorous methods to verify their functional status and, crucially, their capacity for subjective experience. This synthetic substrate path often relies heavily on the development and application of ultra-high-bandwidth, bidirectional Brain-Machine Interfaces (BMIs). Such interfaces serve not merely as input/output channels but as critical tools for gradual integration, functional mapping, and potentially validation – drawing inspiration from proposals for testable machine consciousness, such as those exploring inter-system integration paradigms (as will be further elaborated in Section \ref{sec:masataka-approach}). This route holds the potential for creating substrates with fundamentally greater robustness, engineered fault tolerance, enhanced resilience against environmental hazards, and perhaps even capabilities beyond biological limits. However, this path faces the ‘hard problem’ of consciousness more directly, depending critically on identifying and successfully implementing the correct physical principles or dynamic properties sufficient for instantiating consciousness in a synthetic medium. Success hinges on significant breakthroughs in substrate engineering, ultra-high-fidelity BMI technology capable of seamless, non-disruptive integration, and the development of reliable methods for verifying conscious presence beyond mere functional mimicry. Furthermore, ensuring the continuity of personal identity during any transition or integration process involving a fundamentally different substrate presents unique and formidable theoretical and technical hurdles.

  Pursuing both paths simultaneously provides crucial strategic hedging and risk mitigation. Should the engineering of verifiable consciousness in synthetic substrates prove unexpectedly intractable, or if current assumptions about the sufficiency of certain physical dynamics (e.g., specific computational architectures) turn out to be incorrect, the biologically grounded path offers an alternative route towards significantly extended persistence and resilience. Conversely, if the inherent limitations or vulnerabilities of biological or bio-hybrid systems ultimately prove insurmountable for achieving indefinite continuity or sufficient resilience against catastrophic failure, advancements along the synthetic substrate path offer a potential long-term solution. This duality aligns directly with the core Synconetics principle of seeking robust, engineered solutions while honestly acknowledging current scientific unknowns and technological limitations.

  Crucially, these two paths are not entirely independent; significant synergies exist, and they may eventually converge. Advancements in the sophisticated BCIs required for the later stages of ECP (such as embedding within rich virtual environments or enabling enhanced cognitive control) are direct precursors to the ultra-high-fidelity interfaces essential for the synthetic substrate path. Conversely, insights gained from attempting to engineer Synthetic Consciousness Substrates (SCS)—particularly regarding the minimal dynamic complexity or specific organisational principles required—can directly inform the design criteria and functional targets for the bio-hybrid grafts used in ECP. It is conceivable that biologically grounded approaches like ECP could serve as a vital transitional phase, creating a stabilised, enhanced biological or bio-hybrid platform from which safer, more gradual, and verifiable integration with future synthetic systems might be achieved.

  The deliberate engineering focus of Synconetics enhances the practical feasibility of this roadmap. The ECP path, with its clearly defined intermediate goals in advanced tissue engineering, regenerative medicine for neurological conditions, and improved BCIs, offers tangible near-term therapeutic and potentially commercial value. This creates opportunities for phased, sustainable funding streams, aligning research with demonstrable benefits. The synthetic substrate path, while perhaps representing a longer-term endeavour, involves fundamental research in neuroscience, materials science, physics, and BMI technology that is attractive to governmental and foundational research funding agencies. Its emphasis on developing testable hypotheses and verifiable outcomes, even if focused initially on intermediate measures of complex dynamics or information integration rather than subjective report, makes it more tractable than purely speculative or philosophical approaches to artificial consciousness. Both paths strategically avoid reliance on unproven fundamental physics or distant science-fiction concepts like atomically precise nanotechnology, focusing instead on integrating and aggressively advancing existing technological frontiers in bioengineering, neurotechnology, and complex systems engineering. Nonetheless, the timelines for achieving the ultimate goals of either path remain highly ambitious, and securing consistent, long-term funding—particularly for the more fundamental aspects of the synthetic substrate research—will undoubtedly be challenging and requires demonstrating consistent, verifiable progress against defined milestones.

  % ! MARK: Socio-Economic and Ethical Implications of Realised Synconetics
  \section{Socio-Economic and Ethical Implications}
  \label{sec:economics}

  The potential success of Synconetics methodologies, even within the challenging timeframes we acknowledge, necessitates a departure from purely technical discourse or distant philosophical speculation. If conscious continuity can be reliably engineered, enabling individuals to persist beyond the limitations of their original biological substrate, it precipitates profound socio-economic, political, and ethical questions demanding pragmatic analysis today. The assertion that Synconetics offers a potentially near-term engineering pathway, distinct from indefinite postponement pending future scientific revolutions, compels us to confront these implications not as hypothetical scenarios, but as foreseeable consequences requiring immediate, serious consideration alongside technical research and development.

  The emergence of individuals whose consciousness persists via engineered substrates—whether advanced bio-hybrids or entirely synthetic systems—fundamentally challenges existing legal and political frameworks, which are entirely unprepared for non-biological personhood. How is legal identity defined for an entity potentially lacking a conventional biological body? Questions of citizenship, property ownership, voting rights, and the very basis of legal standing become acutely problematic. Establishing internationally recognised standards for the personhood, rights (such as substrate autonomy, freedom from non-consensual modification, access to environments) and responsibilities (taxation, legal liability) of Synconetic entities represents a monumental political and philosophical undertaking, fraught with potential for inequality and novel forms of exploitation if not proactively addressed.

  Economically, the advent of potentially vastly long-lived or effectively immortal conscious entities promises radical disruption. Can such entities participate meaningfully in labour markets, particularly alongside accelerating AI automation? Assessing their potential for cognitive, creative, or virtual value creation is complex; their existence may necessitate fundamental shifts in economic models, potentially reinforcing arguments for systems like Universal Basic Income if traditional biological labour diminishes further. Furthermore, the significant, ongoing resource demands—energy, computation, physical security, specialised maintenance—for sustaining consciousness-supporting substrates raise critical questions of allocation. What economic models (e.g., subscription, public utility, private ownership) govern access and upkeep, and how can unprecedented societal stratification between those who can afford continuity and those who cannot be avoided? The potential for cost to exacerbate existing inequalities demands careful forethought.

  The infrastructural and logistical realities of supporting a population of Synconetic entities are equally daunting, involving engineering challenges often vastly underestimated. Robust, secure physical and digital infrastructure is paramount. Where are consciousness-supporting substrates housed? What levels of physical security, redundancy against technical failure or environmental catastrophe, and resilience against malicious attack are achievable and sustainable? Centralised hosting creates single points of failure and control, whilst distributed models present immense logistical hurdles. Provider viability is another critical concern: what happens if a commercial or state entity responsible for hosting becomes insolvent, politically unstable, or technologically obsolete? Without clear standards and protocols guaranteeing substrate or data portability—enabling transfer between providers or substrate types without violating continuity—individuals face extreme vulnerability and vendor lock-in. The sheer energy and computational load, especially if entities interact within rich virtual environments, also poses significant questions about long-term global sustainability.

  Perhaps most profoundly, the successful realisation of Synconetics challenges fundamental societal notions of life, death, identity, and community. How will society perceive these entities—as ‘alive’, ‘post-biological’, or something entirely new? How do existing relationships, inheritance laws, legacy considerations, and social security systems adapt? The psychological well-being of individuals undergoing transition and potentially existing indefinitely, perhaps within environments vastly different from baseline biological reality, presents significant risks and necessitates novel forms of support. Maintaining existential meaning under such conditions is a critical, open question. Ensuring equitable access and mitigating the potential for coercion (e.g., societal pressure to transition) are paramount ethical considerations. Finally, defining ‘death’ for a Synconetic entity and establishing ethical end-of-life protocols—managing substrate failure, irreversible cognitive decline, or respecting an individual’s voluntary wish to cease existence—represents entirely uncharted territory demanding sensitive, cross-disciplinary deliberation.

  The potential near-term feasibility advocated by Synconetics thus transforms these issues from speculative fiction into urgent matters for contemporary policy, ethics, and engineering. The stark contrast between this potential and the current lack of serious planning underscores the imperative for proactive engagement. Addressing the legal, economic, infrastructural, and ethical dimensions cannot be postponed; it must occur in parallel with technical research and development. This proactive, transdisciplinary effort is essential to mitigate the risks of societal disruption, inequality, and catastrophic failure, ensuring that the pursuit of engineered continuity aligns with broadly shared human values. It is a core tenet of the Synconetics approach that responsible engineering necessitates foresight into its societal consequences.

  % ! MARK: Conclusion and Call to Action
  \section{Conclusion and Call to Action}
  \label{sec:conclusion}
  This essay has introduced Synconetics, a scientific and engineering discipline founded upon the conviction that biological death, understood fundamentally as a failure of the consciousness-supporting substrate, represents a tractable engineering challenge. We have argued that by rigorously prioritising the uninterrupted physical continuity of the processes underpinning individual consciousness—the preservation of the unique 4D process-world-line—and concentrating on tangible, buildable systems grounded in established science, Synconetics charts a more robust, ethically defensible, and ultimately achievable course than paradigms reliant on destructive replication or unverified philosophical assumptions, such as strong computationalism. Its framework offers a pragmatic pathway towards ensuring the persistence of conscious existence.

  The methodologies currently being developed within the Synconetics framework, exemplified by the Ectopic Cognitive Preservation strategy (Section \ref{sec:daniel-approach}) focused on gradual bio-hybrid replacement, and complemented by approaches centred on advanced Brain-Machine Interface integration for probing and potentially validating synthetic substrates (as anticipated in Section \ref{sec:masataka-approach}), serve as initial, concrete demonstrations of this potential. They affirm that research and development aligned with Synconetics principles can commence immediately, strategically leveraging existing and foreseeable advancements across neuroscience, bioengineering, materials science, and related fields. This potential feasibility, suggesting meaningful progress within decades rather than indefinite centuries, transforms the profound socio-economic, legal, and ethical questions accompanying engineered conscious continuity from distant speculations into urgent matters demanding immediate, serious consideration (Section \ref{sec:economics}). Proactive, transdisciplinary planning and societal dialogue are not merely advisable; they are imperative to navigate the immense societal shifts this technology could precipitate.

  We contend, therefore, that a significant redirection of focus and resources is necessary within the broader constellation of research aiming to overcome biological limitations. A paradigm shift is required, moving decisively towards the direct engineering of continuity and substrate resilience. This involves embracing the complexities of physical instantiation and continuous process dynamics, rather than pursuing potentially flawed or existentially risky shortcuts predicated on abstract information patterns or destructive scanning alone.

  The advancement of Synconetics, however, cannot be the work of isolated groups; it demands a concerted, collaborative, and deeply transdisciplinary effort. We issue a call to action to researchers, engineers, clinicians, ethicists, policymakers, entrepreneurs, and funders worldwide to engage actively with this nascent field. We invite scientists and engineers to \textit{advance the research frontier} by pursuing fundamental research and targeted engineering development aligned with Synconetics principles; opportunities exist for postgraduate research exploring consciousness mechanisms and BMI integration (e.g., with Prof. Watanabe’s group at the University of Tokyo) and for applied RnD within dedicated ventures. We encourage engagement with, and support for, organisations \textit{translating Synconetics principles into practice}, such as Eightsix Science (currently seeking technical collaborators, funding, and grant support for its ECP approach). We urge innovators to \textit{foster diversity and progress} by launching new research projects or companies exploring alternative continuity-preserving strategies; a healthy ecosystem of complementary approaches will strengthen the entire field.

  Furthermore, we call upon the community to \textit{engage in critical dialogue}: connect with the authors and other researchers to discuss these concepts, rigorously challenge assumptions, and collaboratively refine the Synconetics framework. Join the nascent community discussions (e.g., via the established Discord server) to share insights and foster collaboration. Help \textit{disseminate and develop knowledge} by sharing this whitepaper and engaging peers in substantive discussion. Contribute to future knowledge-building efforts, such as the planned comprehensive book, \textit{Synthetic Consciousness}; we actively seek co-authors from diverse disciplines, particularly medicine, law, economics, and ethics, to ensure a truly comprehensive and transdisciplinary perspective. Finally, support or participate in initiatives designed to \textit{convene the community}, potentially including a dedicated Synthetic Consciousness Conference, to consolidate research findings and catalyse interdisciplinary exchange.

  Synconetics represents more than a theoretical exercise or a distant dream; it is a call to apply the full power of rigorous engineering principles, tempered by ethical foresight, to one of humanity’s oldest and most profound challenges. By maintaining an unwavering focus on verifiable physical continuity and the development of buildable, reliable systems, we can begin to move beyond speculative fiction towards tangible progress in ensuring the persistence of human consciousness. We welcome all who share this vision and commitment to join us in building this critical field.

  % ! ========================
  % ! # MARK: References, etc.
  % ! ========================

  \pagebreak
  \bibliographystyle{../../templates/custom-apa}
  \bibliography{references/bibliography}
  \nocite{*}

\end{sloppypar}
\end{document}